\section{The store machine}\label{sec:machine}

A record is kept in a \emph{store}: a content (objects at addresses and relations between addresses), a guard register (predicates the
content must satisfy), a list of snapshots, and a ledger.  A store is \emph{admissible} if every relation uses stored addresses and every
guard holds.

\paragraph{Instructions (MML).}  \textsc{store}\((a,x)\) puts an object at an address, \textsc{link}\((a,b)\) adds a relation,
\textsc{snapshot} saves the content, \textsc{rollback}\((k)\) restores the \(k\)-th snapshot, \textsc{scrub} is a no-op used for cleaning
passes, and \textsc{export}\((a)\) releases a stored object.

\paragraph{Transition semantics (MMA).}  Each instruction has a raw effect, which may be undefined (a rollback to a missing snapshot, an
export of a missing address).  A step executes the raw effect and checks the guards on the result:
\[
\mathrm{exec}(s,i)=\begin{cases}
(s\ \text{with entry } (i,\textsc{reject}),\ \textsc{reject}) & \text{if the raw effect is undefined},\\
(s'\ \text{with entry } (i,\textsc{cont}),\ \textsc{cont}) & \text{if it gives } s' \text{ and } s' \text{ is admissible},\\
(s\ \text{with entry } (i,\textsc{guard}),\ \textsc{guard}) & \text{otherwise.}
\end{cases}
\]
A result that violates a guard is never committed; the old content stays and the ledger records the defect.  A program is a list of
instructions executed in order.

\paragraph{Certificates.}  The record distinguishes a \emph{candidate}, which the machine can produce for any statement, from a
\emph{certificate}, which contains a proof.  A candidate becomes a certificate only through a proof supplied from outside the machine.  No
instruction produces a certificate.
